\subsection{Quadrature lemmas and the dual formulation of the entropy}\label{sec:D01}
From Section~\ref{sec:perfil} we use the bounds $c/n\le\pi\le M/n$, the fact that the diagonal vacancies are
at least $\eta$, and the estimate $\|z^*-z^0\|_\infty=O(n^{-1})$ in the row and column gauge.
The continuous profile $\nu$ and its positive factors are Lipschitz.
The estimates hold for all sufficiently large $n$, with constants independent of $n$.

\begin{lema}[Midpoint quadrature with a small remainder]\label{lem:puntosmedios}
For a fixed Lipschitz function $f$ on $[0,1]^d$, $d=1,2$,
the uniform midpoint rule satisfies
$Q_n(f)-\int f=o(n^{-1})$.
\end{lema}
\begin{proof}
Let $h=1/n$, and let $c_n(x)$ be the center of the cell containing $x$.
Since Lipschitz functions are differentiable almost everywhere (Rademacher's theorem),
\[
 f(c_n(x))-f(x)=\nabla f(x)\cdot(c_n(x)-x)+r_n(x),
 \qquad r_n(x)/h\longrightarrow0\quad\text{almost everywhere}.
\]
The quotient $|r_n|/h$ is bounded by a constant.
The field $w_n=(c_n-x)/h$ is uniformly bounded and has
mean zero on each cell. For every continuous field $v$,
$\int v\cdot w_n\to0$ by uniform continuity.
Approximating $\nabla f\in L^1$ by continuous fields proves the same assertion
for $\nabla f$. Integrating the identity and applying dominated convergence
to the remainder gives the result.
The proof does not invoke pointwise evaluation of arbitrary $W^{1,1}$ functions
in two dimensions.
\end{proof}

\begin{lema}[One-dimensional interior grid]
For $f$ Lipschitz on $[a,b]$, with the endpoints on the grid of spacing $1/n$,
\[
 \sum_{a<k/n<b}f(k/n)
 =n\int_a^b f-\frac{f(a)+f(b)}2+o(1).             \tag{D1.1}
\]
\end{lema}
\begin{proof}
On each interval of length $h$, the trapezoidal error is
$\int (x-\text{center})f'(x)\,dx$.
After division by $h$, its kernel is bounded and has mean zero
on each interval. Approximating $f'\in L^1$ by continuous functions
gives a total error of $o(h)$. Solving for the interior sum proves (D1.1).
\end{proof}

\paragraph{The exact dual functional.}
We use the operators in (TR3), $\mathcal Tz=\sum_{L\ni c}z_L$ and
$(\mathcal R_fq)_L=n^{-1}\sum_{c\in L}q_c$.
Let $\chi_n$ be the inverse of $\psi_n(q)=q(1-q/n)^3$ on the small branch, and set
$q(z)=\chi_n(e^{\mathcal Tz})$ and
\[
 A_n(q)=4q+3n\log(1-q/n),
 \quad E_n(z)=\sum_{D,A}e^{z_L}-\sum_{\text{all }L}z_L
                         +\frac1n\sum_c A_n(q(z)_c).       \tag{D1.2}
\]
Since $dA_n/d\log\psi_n=q$, the Euclidean gradient of $E_n$ is exactly
\[
 F_n(z)=(\mathcal R_Rq-1,\mathcal R_Cq-1,
                 e^{z_D}+\mathcal R_Dq-1,e^{z_A}+\mathcal R_Aq-1).
\]
Section~\ref{sec:P05} proves $F_n(z^*)=0$ and
$\|F_n(z^0)\|_\infty=O(n^{-1})$. On the segment joining the two points,
$DF_n=d+\mathcal R(w\mathcal T)$, with $0\le d\le1$ and
$w\le C_w$ (Section~\ref{sec:P06}, (TP1)). Each line contains at most $n$ squares
and each square lies on four lines, so
$\|\nabla^2E_n\|_{\infty\to\infty}\le1+4C_w$.
Since there are $6n-2$ variables, Taylor's formula with integral remainder gives
\[
 E_n(z^*)-E_n(z^0)=O(n^{-1}).                         \tag{D1.3}
\]
Indeed, the linear term is bounded by
$(6n-2)\|F_n(z^0)\|_\infty\|z^*-z^0\|_\infty$,
and the quadratic term by a constant times
$(6n-2)\|z^*-z^0\|_\infty^2$.

\paragraph{The entropy at the critical point.}
At the maximizer, $q^*=n\pi$, the row and column sums equal one,
$\sum_{D,A}h_L=2n-2$, and $\sum_c q_c^*/n=n$.
Stationarity implies
\[
 \log\pi_c+3\log(1-\pi_c)
       =z^*_{R}+z^*_{C}+\log h_D+\log h_A-\log n.
\]
Upon substitution into the Bethe entropy, the terms for each diagonal
combine by $m_L+h_L=1$, giving
\[
 H(\pi)=n\log n-\sum_Lz_L^*+3\sum_c\log(1-\pi_c)
       =n\log n+E_n(z^*)-(6n-2).                    \tag{D1.4}
\]
The term $n\log n$ comes from the exact discrete mass, before any quadrature.
